\section*{Methods}
The mathematics was produced in the order: computation, formalization, prose. Candidate statements were explored in a computer algebra system. Selected load-bearing statements were fixed by the author and frozen in separate files at designated milestones, before the corresponding production proofs or subsequent refactorings were accepted; the frozen layer is diffed against the production tree at every build. Computer-algebra and finite-instance checks — positive, negative and near-miss instances against an independent reimplementation — were used as adversarial validation of algebraic identities and translations. They do not certify that a formal statement models the intended problem: semantic correspondence between the intended mathematical problem, the paper statement and the Lean type remained subject to author review, and two earlier misformalizations of that kind were caught by review, not by the kernel. Tactic-level proof search in Lean was delegated to a large language model acting as a coding agent, in single-threaded sessions with no nested delegation; the author independently executed the clean build, the theorem audit and the axiom-footprint reports on the final source tree, model-generated proof terms were accepted only after kernel checking and reconciliation with the human-readable proof, and no agent self-report was accepted as evidence. A second model, from a different vendor, served as an adversarial reader of statements and drafts; its objections were recorded with dispositions, and those that were correct led to corrections incorporated here. Randomised and exhaustive numerical checks were used to look for counterexamples and to corroborate new lemmas before they were written down; they are evidence of the absence of certain errors and are never treated as proof.

